% Native reconstruction of French original PDF268/printed263.
\clearpage
\section*{本讲目录}
\begingroup\small
\noindent\hfill 原书分讲页码\par\smallskip
\begin{tabular}{@{}p{.85\linewidth}r@{}}
\multicolumn{2}{c}{第一章\quad 三角范畴}\\[5pt]
\S1. \hyperref[VIII:1]{定义与例子} & \\
\quad 1. \hyperref[VIII:1-1]{三角范畴的定义. 初步性质} & 1\\
\quad 2. \hyperref[VIII:1-2]{三角范畴的例子} & 5\\
\quad 3. \hyperref[VIII:1-3]{正合函子和上同调函子的例子} & 10\\[4pt]
\S2. \hyperref[VIII:2]{商范畴} & \\
\quad 1. \hyperref[VIII:2-1]{厚子范畴. 乘法系} & 13\\
\quad 2. \hyperref[VIII:2-2]{等价的泛问题} & 15\\
\quad 3. \hyperref[VIII:2-3]{分式演算} & 16\\
\quad 4. \hyperref[VIII:2-4]{商范畴的性质} & 19\\
\quad 5. \hyperref[VIII:2-5]{商函子的性质} & 21\\
\quad 6. \hyperref[VIII:2-6]{正交性} & 24\\[7pt]
\multicolumn{2}{c}{第二章\quad 对阿贝尔范畴的应用}\\[5pt]
\S1. \hyperref[VIII:3]{导出范畴} & \\
\quad 1. \hyperref[VIII:3-1]{阿贝尔范畴的导出范畴} & 28\\
\quad 2. \hyperref[VIII:3-2]{$\ext$ 的研究} & 34\\[4pt]
\S2. \hyperref[VIII:4]{导出函子} & \\
\quad 1. \hyperref[VIII:4-1]{导出函子的定义} & 37\\
\quad 2. \hyperref[VIII:4-2]{导出函子的存在性} & 40\\
\quad 3. \hyperref[VIII:4-3]{杯积. 复合} & 44\\[4pt]
\S3. \hyperref[VIII:6]{例子} & \\
\quad 1. \hyperref[VIII:6-1]{函子 $\underline{\operatorname{Hom}}$} & 46\\
\quad 2. \hyperref[VIII:6-2]{层的 $\Tor$} & 47\\
\quad 3. \hyperref[VIII:6-3]{逆像的左导出函子} & 48
\end{tabular}
\par\medskip
\noindent 本目录保留原书标题与分讲页码; 原书部分目录标题与正文标题略有不同. 本译本正文页码见卷首目录.
\endgroup
\clearpage
